\section{Related Work}
\label{sec:related}

\para{Literature-grounded ideation.} A large body of work grounds idea generation
in retrieved prior art. \scimon{} retrieves inspirations and runs an iterative
``novelty-boosting'' loop that forces each idea to differ from its $k$ nearest
neighbors, raising human-judged novelty for $46$--$58\%$ of
ideas~\citep{wang2023scimon}. ResearchAgent traverses citation graphs and an
entity knowledge store~\citep{baek2024researchagent}; Chain-of-Ideas induces the
developmental trend of a field before proposing~\citep{li2024coi}; Nova plans and
retrieves each round to keep ideas novel~\citep{hu2024nova}; and MOOSE-Chem composes
hypotheses from background plus deliberately non-neighboring
inspirations~\citep{yang2024moosechem}. We instantiate this family as our literature-
comparison regime, but---to keep the ablation leakage-controlled and single-model---we
differentiate against the model's own accumulating idea pool rather than an external
corpus. Unlike these systems, our goal is not to build a better literature-grounded
generator but to place literature comparison \emph{on the same footing} as other
grounding mechanisms.

\para{Data-driven and empirical grounding.} \hypogenic{} induces explanatory
hypotheses from labeled examples and refines them with a bandit-style
reward tied to predictive accuracy, beating few-shot prompting by $3.3$--$31.7\%$ and
even fine-tuned oracles~\citep{zhou2024hypogenic}. Literature Meets Data shows
literature and data are complementary, improving \OOD{} accuracy by
$8.97\%$ over few-shot~\citep{liu2024literature}. SGA closes the loop with a
differentiable simulator, where removing the simulation grounding collapses
performance by orders of magnitude~\citep{ma2024sga}. \hypobench{} benchmarks these
methods at scale and finds data-driven and literature+data methods dominate
inference-only and zero-shot-generation baselines, especially
\OOD~\citep{liu2025hypobench}. Our induction and empirical regimes are direct,
minimal instantiations of this line; our contribution is to measure them against the
\emph{same} ungrounded control and each other, and to report when data grounding
\emph{overfits} rather than transfers.

\para{Proposal-writing and structured ideation.} Several systems ground ideation by
writing full structured proposals (problem $\rightarrow$ method $\rightarrow$
experiments) that are critiqued and revised~\citep{baek2024researchagent,si2024llm}.
\citet{si2024llm} find LLM proposals are judged more novel than expert humans but
slightly less feasible. We include proposal writing as a regime and find that, when
distilled into decision rules, it can \emph{reduce} usefulness and diversity---a
concrete case of an idea looking thorough while being less useful.

\para{Evaluating ideas.} Because ``is this novel?'' prompts are unreliable, recent
work grounds the \emph{judgment} in retrieved literature: RAG-Novelty and
literature-grounded assessment beat ungrounded LLM-as-judge
prompting~\citep{lin2024schnovel,shahid2025litnovelty}, and idea-level benchmarks
catalog the trade-offs~\citep{guo2024ideabench,liu2025researchbench}. The consistent
caveat---\emph{novelty $\neq$ value}, sharpened by the ideation--execution flip of
\citet{si2024llm}---motivates our choice to score grounded usefulness (downstream
\IND/\OOD{} accuracy) rather than a bare novelty score, and to test the
novelty--usefulness relationship directly.

\para{Our position.} Every method above grounds in a single way and beats an
ungrounded prompt measured with a weaker model. We hold task, model, generation
budget, and evaluation fixed and ablate the grounding \emph{mechanisms} against one
another under one protocol. To our knowledge this head-to-head comparison does not
exist, and it yields a partly negative, mechanism-specific picture that the
single-grounding literature could not surface.
